\documentclass[10pt]{article}
\usepackage[UTF8]{ctex}
\usepackage{amsmath, amssymb}
\usepackage{geometry}
\geometry{a4paper, margin=1.5cm}
\usepackage{booktabs}

\title{数学建模应用与探究题参考答案}
\author{}
\date{}

\begin{document}
\maketitle

\section{得分波动分析参考答案}

在试卷结构权重 $w_1=0.2,\ w_2=0.6,\ w_3=0.2$ 下，学生甲（中等生）取 $p_1=0.9,\ p_2=0.7,\ p_3=0.2$，学生乙（学霸）取 $p_1=1.0,\ p_2=0.9,\ p_3=0.8$。代入期望公式 $E(R)=\sum w_i p_i$ 与方差公式 $D(R)\approx\sum w_i p_i(1-p_i)$，计算得：
$$
E_{\text{甲}}=0.2\times0.9+0.6\times0.7+0.2\times0.2=0.18+0.42+0.04=0.6400,
$$
$$
D_{\text{甲}}=0.2\times0.9\times0.1+0.6\times0.7\times0.3+0.2\times0.2\times0.8=0.018+0.126+0.032=0.1760;
$$
$$
E_{\text{乙}}=0.2\times1.0+0.6\times0.9+0.2\times0.8=0.2+0.54+0.16=0.9000,
$$
$$
D_{\text{乙}}=0.2\times1.0\times0+0.6\times0.9\times0.1+0.2\times0.8\times0.2=0+0.054+0.032=0.0860.
$$
可见甲的期望得分率显著低于乙，且甲的方差远大于乙，说明中等生成绩波动明显大于学霸。

从函数结构看，方差的核心项是 $p(1-p)$，该函数在 $p=0.5$ 处取得最大值 $0.25$，在 $p\to0$ 或 $p\to1$ 时趋近于 $0$。中等生的得分率 $p_2=0.7$、$p_3=0.2$ 均处于 $0.2\sim0.8$ 区间，对应的 $p(1-p)$ 值较大；而学霸的 $p_i$ 均接近 $1$，$p(1-p)$ 值很小。此外，中档题权重 $w_2=0.6$ 最高，进一步放大了中等生在 $p_2=0.7$ 处的方差贡献。因此，中等生波动更大具有统计上的必然性。

\section{估分偏差分析参考答案}

已知 $E(Y)=0.2r^2+0.7r+0.1$，系统偏差函数 $b(r)=E(Y)-r=0.2r^2-0.3r+0.1$。解方程 $b(r)=0$：
$$
0.2r^2-0.3r+0.1=0 \;\Longrightarrow\; 2r^2-3r+1=0 \;\Longrightarrow\; (2r-1)(r-1)=0,
$$
得两个零点 $r=0.5$ 与 $r=1.0$。这意味着当自评得分率为 $0.5$ 或 $1.0$ 时，期望意义下估分无系统偏差。由二次函数开口向上可知：当 $r<0.5$ 时 $b(r)>0$，即低分段学生整体倾向高估成绩；当 $0.5<r<1$ 时 $b(r)<0$，即中高分段学生整体倾向低估成绩；而 $r=1.0$（顶尖生）估分与实际期望完全一致。

方差函数 $D(Y)=0.09-0.08r-0.01r^2$。求导得 $D'(Y)=-0.08-0.02r<0$（$r\in[0,1]$），故 $D(Y)$ 严格单调递减。这表明自评得分率越高（即水平越扎实），实际得分的随机波动越小，发挥越稳定，这与“掌握越扎实，考试发挥越稳定”的经验完全吻合。

取 $r=0.60$ 与 $r=0.90$ 分别代表中等生和学霸，计算得：
$$
b(0.6)=0.2\times0.36-0.3\times0.6+0.1=0.072-0.18+0.1=-0.0080,
$$
$$
D(0.6)=0.09-0.08\times0.6-0.01\times0.36=0.09-0.048-0.0036=0.0384,
$$
$$
\text{MSE}(0.6)=0.0384+(-0.0080)^2=0.0384+0.000064=0.0385;
$$
$$
b(0.9)=0.2\times0.81-0.3\times0.9+0.1=0.162-0.27+0.1=-0.0080,
$$
$$
D(0.9)=0.09-0.08\times0.9-0.01\times0.81=0.09-0.072-0.0081=0.0099,
$$
$$
\text{MSE}(0.9)=0.0099+(-0.0080)^2=0.0099+0.000064=0.0100.
$$
可见两者系统偏差 $b(r)$ 相同，但学霸的方差 $D(Y)$ 远小于中等生，因此均方误差 MSE 也更小，即估分更准。MSE 同时考虑了系统偏差和随机波动，能够更全面地评价估分精度，避免了只看偏差或只看方差的片面性。

\section{学科时间分配模型参考答案}

给定每周总时间 $T=30$ 小时，各科参数如题中表格所示。首先计算每科“维持不下滑”的最少投入时间，公式为
$$
t_{\min}=\frac{\theta(S_{\text{cur}}-S_0)}{\eta(M-S_{\text{cur}})}.
$$
代入数据得：

\begin{center}
\begin{tabular}{lcccccc}
\toprule
科目 & 语文 & 数学 & 英语 & 历史 & 政治 & 地理 \\
\midrule
$t_{\min}$（小时/周） & 3.33 & 2.67 & 2.50 & 2.67 & 1.67 & 1.67 \\
\bottomrule
\end{tabular}
\end{center}

计算过程示例（数学）：$t_{\min}=\frac{0.20\times(90-70)}{0.05\times(120-90)}=\frac{4}{1.5}\approx2.67$。求和得总保底时间 $T_{\min}=14.50$ 小时，剩余可突破时间 $\Delta T=30-14.50=15.50$ 小时。

若将 $\Delta T$ 全部投入数学，则 $t_{\text{数}}=2.67+15.50=18.17$，代入稳态分数公式：
$$
S_{\text{数}}^*(18.17)=\frac{0.05\times18.17\times120+0.20\times70}{0.05\times18.17+0.20}=\frac{109.02+14}{0.9085+0.20}=\frac{123.02}{1.1085}\approx110.98,
$$
提升量 $\Delta S=110.98-90=20.98$ 分。若全部投入英语，$t_{\text{英}}=2.50+15.50=18.00$，
$$
S_{\text{英}}^*(18.00)=\frac{0.02\times18.00\times125+0.10\times95}{0.02\times18.00+0.10}=\frac{45+9.5}{0.36+0.10}=\frac{54.5}{0.46}\approx118.48,
$$
提升量 $\Delta S=118.48-105=13.48$ 分。单投数学提分更多，推荐优先突破数学。

若同时突破数学和英语，设数学分配 $x$ 小时（$0\le x\le 15.50$），英语分配 $15.50-x$ 小时。总稳态分数为
$$
S_{\text{总}}(x)=\bigl(\sum_{i\neq\text{数,英}}S_{i,\text{cur}}\bigr)+S_{\text{数}}^*(2.67+x)+S_{\text{英}}^*(2.50+15.50-x).
$$
由等边际收益原理，最优 $x$ 满足
$$
\left.\frac{dS_{\text{数}}^*}{dt}\right|_{t=2.67+x}=\left.\frac{dS_{\text{英}}^*}{dt}\right|_{t=18.00-x},
$$
其中 $\frac{dS^*}{dt}=\frac{\eta\theta(M-S_0)}{(\eta t+\theta)^2}$。代入参数可解得最优 $x\approx9.23$ 小时（数学），英语相应为 $6.27$ 小时。此时数学稳态分约为 $105.68$，英语稳态分约为 $117.85$，总提升约为 $26.53$ 分，高于单科突破方案。

\subsection*{拓展：直接求解约束优化问题}
若不采用两阶段策略，而是直接求解 $\max\sum_i S_i^*(t_i)$ 满足 $\sum t_i=30,\ t_i\ge0$，计算得到的全局最优时间分配及对应稳态分数如下：

\begin{center}
\begin{tabular}{lcccccc}
\toprule
科目 & 语文 & 数学 & 英语 & 历史 & 政治 & 地理 \\
\midrule
最优时间（小时/周） & 5.10 & 8.78 & 6.07 & 3.93 & 3.06 & 3.06 \\
对应稳态分 & 107.63 & 104.35 & 111.45 & 81.92 & 77.18 & 77.18 \\
\bottomrule
\end{tabular}
\end{center}
六科总稳态分为 $559.70$，相对当前总分（$105+90+105+80+75+75=530$）提升 $29.70$ 分。该结果说明，全局优化比仅在两科之间分配剩余时间能获得更多收益。

\section{劳逸结合最优配置参考答案}

学习产出函数为 $P(t_1,t_2,t_3)=(4+0.9t_1)(1+0.2t_2)(1+0.15t_3)$，约束 $t_1+t_2\le3,\ 0\le t_3\le1,\ t_i\ge0$。

极端刷题策略 $(t_1,t_2,t_3)=(3,0,0)$ 时，
$$
P_1=(4+0.9\times3)\times(1+0)\times(1+0)=(4+2.7)=6.7000.
$$
平衡策略 $(t_1,t_2,t_3)=(1.5,1.5,1)$ 时，
$$
P_2=(4+0.9\times1.5)\times(1+0.2\times1.5)\times(1+0.15\times1)=(4+1.35)\times(1+0.3)\times1.15=5.35\times1.3\times1.15=7.9982.
$$
$P_2 > P_1$，说明加入深度思考与运动后总产出显著提升。这是由于乘法结构使得各因子相互协同：刷题提供知识积累基础，思考提升思维锐度，运动增强精力水平，三者乘积放大了总效果，避免了单一因子成为短板。

由于 $P$ 关于 $t_3$ 单调递增，取 $t_3=1$。由 $t_1+t_2=3$ 得 $t_1=3-t_2$，代入得
$$
P(t_2)=(4+0.9(3-t_2))\times(1+0.2t_2)\times1.15=(4+2.7-0.9t_2)\times(1+0.2t_2)\times1.15=1.15\,(6.7-0.9t_2)(1+0.2t_2).
$$
展开得 $P(t_2)=1.15\,(6.7+1.34t_2-0.9t_2-0.18t_2^2)=1.15\,(6.7+0.44t_2-0.18t_2^2)$。这是关于 $t_2$ 的二次函数，开口向下，顶点在 $t_2^*=\frac{0.44}{2\times0.18}\approx1.222$ 小时。于是 $t_1^*=3-1.222=1.778$ 小时，$t_3^*=1$ 小时。最大产出 $P_{\max}=1.15\,(6.7+0.44\times1.222-0.18\times1.222^2)=1.15\,(6.7+0.5377-0.268)=1.15\times6.9697\approx8.015$。该最优值高于极端刷题策略，且与平衡策略 $7.998$ 接近，进一步说明适度运动和思考有助于提升学习效率。

\section*{综合结论与使用建议}

以上四题共同揭示了一条主线：有效学习策略不是将所有资源集中于单一环节，而是在概率波动、边际收益、系统偏差和精力协同之间寻求平衡。第一题解释了中等生波动大的统计根源；第二题说明了高水平学生估分更准的原因（方差小且系统偏差可控）；第三题给出了可执行的时间分配优化路径（先保底后突破，或全局优化）；第四题从乘法协同角度证明了劳逸结合的数学必然性。这套题目不仅训练计算能力，更培养将现实问题转化为数学结构并反向解释现实的建模思维。

\end{document}